% Píše: aplikace. Sekce (\section) o testování aplikace; kapitolu
% "Testování a ověření" otevírá Chapters/testovani.tex.
\section{Testování mobilní aplikace}
\label{sec:testovani-aplikace}

Aplikace se testuje na třech úrovních: automatickými testy na počítači, ručně v~telefonu proti
simulovanému domu a~se skutečným Raspberry Pi. Nová funkce i~oprava chyby začínala testem,
který bez změny kódu selhal, a~teprve potom vznikl kód, se kterým test prošel.

\subsection{Automatické testy}

Automatické testy běží na počítači v~prostředí Java Virtual Machine (JVM), bez telefonu
i~emulátoru, a~spouští je příkaz \kod{./gradlew testDebugUnitTest}. Obrazovky v~nich vykresluje
knihovna Robolectric~\cite{app-robolectric}, která na JVM napodobuje Android. Ovládá je testovací
rozhraní Compose~\cite{app-compose-testy}: najde prvek podle textu nebo popisu pro čtečku
obrazovky, klepne na něj a~ověří, co se stalo. Testy ViewModelu běží ve virtuálním čase knihovny
kotlinx-coroutines-test~\cite{app-coroutines-testy}, takže limity v~řádu sekund ověří okamžitě
a~pokaždé stejně. Místo skutečného zdroje dat dostávají falešný zdroj, kterému test nastavuje
Stav a~Nastavení a~který zaznamenává zápisy a~příkazy.

Aplikace má 1\,173~automatických testů a~všechny procházejí. Tabulka~\ref{tab:app-testy}
ukazuje, jak se dělí podle částí kódu.

\begin{table}[H]
\centering
\caption{Automatické testy aplikace podle částí kódu}
\label{tab:app-testy}
\begin{tabular}{|p{0.20\textwidth}|p{0.09\textwidth}|p{0.59\textwidth}|}
\hline
\textbf{Část kódu} & \textbf{Testů} & \textbf{Co testy ověřují} \\ \hline
obrazovky & 457 & texty a~stavy prvků, akce po klepnutí, zašedlé ovladače, rozvržení na telefonu
  a~tabletu, texty při zvětšeném písmu \\ \hline
ViewModel & 252 & zpoždění a~slučování zápisů, potvrzení do 3~s, hlášky po chybě, deník, historie
  grafu \\ \hline
simulace & 167 & pravidla domu podle rozhraní: termostat, automatika světla, úseky dne, odmítnutí
  neplatných zápisů a~příkazů \\ \hline
snímky obrazovek & 119 & vykreslení obrazovek ve všech důležitých stavech do obrázků \\ \hline
doménová logika & 69 & úsek dne, prahy termostatu, texty světla a~chyb, český formát čísel \\ \hline
Bluetooth LE & 46 & odstupy pokusů o~připojení, výběr domu po hledání, rozestup startů hledání,
  kontrola MTU, oprávnění podle verze Androidu \\ \hline
datový model & 39 & převod JSON, tolerance k~neznámým klíčům a~číslům jako \kod{50.0}, celá čísla
  v~zápisech bez desetinné tečky \\ \hline
ostatní & 24 & aktivita a~pomocné nástroje testů \\ \hline
\end{tabular}
\end{table}

Výpis~\ref{lst:app-test} ukazuje test z~části ViewModel. Hlídá, že svítící světlo po výběru
barvy neukáže na okamžik starou barvu (oddíl~\ref{sec:app-data}): vybraná barva musí zůstat,
dokud ji nepotvrdí Nastavení i~Stav. Funkce \kod{runCurrent} spustí všechno, co ve virtuálním
čase čeká na provedení, a~proměnná \kod{fake} je falešný zdroj dat.

\begin{lstlisting}[language=kotlin, breaklines=true, caption={Test potvrzení vybrané barvy světla}, label=lst:app-test]
@Test
fun aShiningLightKeepsTheStoredLookDraftUntilTheStatusShowsItToo() = vmTest {
    val vm = connectedViewModel()
    val red = LightState(on = true, color = listOf(255, 0, 0), brightness = 100, manual = true)
    fake.status.value = statusWithLight(red)
    runCurrent()
    vm.setLightColor(Rgb(0, 0, 255))
    runCurrent()

    fake.settings.value = DEFAULTS.copy(lightColor = listOf(0, 0, 255))
    runCurrent()
    assertNotNull(vm.state.lightDraft)
    assertEquals(red.copy(color = listOf(0, 0, 255)), vm.state.shownLight)

    fake.status.value = statusWithLight(red.copy(color = listOf(0, 0, 255)))
    runCurrent()
    assertNull(vm.state.lightDraft)
}
\end{lstlisting}

Testy snímků obrazovek vykreslí 114~obrázků: obrazovky na telefonu na výšku i~na šířku, na
tabletu, v~tmavém motivu a~s~písmem zvětšeným na 150~\%. Kromě uložení obrázku každý test
ověří, že se žádný text neořízne, nezalomí uprostřed slova a~nepřekrývá s~jiným. Obrázky se
procházejí ručně po každé změně vzhledu a~pocházejí z~nich i~snímky v~této dokumentaci.

Automaticky se neověřuje samotná komunikace s~Bluetooth zásobníkem Androidu, protože zpětná
volání třídy \kod{BluetoothGatt} na JVM nefungují. Testy pokrývají logiku klienta Bluetooth LE,
která na nich nezávisí, a~spojení se skutečným domem se zkouší ručně.

\subsection{Ruční testy}

Ruční zkouška se řídí kontrolním seznamem, který je součástí zdrojových kódů aplikace. Seznam
prochází první spuštění a~oprávnění, všechna tlačítka simulace, zápisy nastavení a~jejich
odmítnutí, ovládání světla, deník, rozvržení na šířku, na tabletu, se zvětšeným písmem
a~v~tmavém motivu a~obnovu spojení po vypnutí Bluetooth, vzdálení z~dosahu a~restartu
Raspberry Pi. U~každého kroku uvádí, co přesně se má na obrazovce objevit. Aplikace se zkouší
na telefonu Samsung Galaxy A25 s~Androidem~16 a~na tabletu Samsung Galaxy Tab A8
s~Androidem~14.

\subsection{Test se skutečným Raspberry Pi}

První společný test se skutečným Raspberry Pi ověřil, že se aplikace připojí s~MTU 517,
dostává notifikace Stavu každou sekundu a~zobrazuje živá data. Test ukázal tři odchylky na
straně domu: vlhkost chodila jako 0, i~když čidlo BMP280 vlhkost neměří a~podle rozhraní má
chodit \kod{null}, osvětlení chodilo jako \kod{null}, protože program počítal s~jiným čidlem,
a~hodiny Raspberry Pi šly asi o~2~minuty napřed. Zjištění se předala straně backendu, která
první dvě opravila (oddíl~\ref{sec:problemy-hw}). Rozdíl hodin aplikace ukazuje a~srovná ho
příkazem \kod{cas}. Po testu backend dodal závazný popis rozhraní (kapitola~\ref{chap:rozhrani})
a~aplikace se mu přizpůsobila, například v~názvu klíče pro rozsvícení při pohybu a~v~kódech
poruch čidel.
